\documentclass{icmmcm}
\usepackage{url}              % For formatting URLs and other web or
                              % file references.
\usepackage{mflogo}           % Provides the METAFONT logo.
\usepackage[dvips]{graphics}  % For importing graphics for use with
                              % dvips.
%\usepackage[pdftex]{graphics} % For importing graphics for use with
                              % pdftex or pdflatex.
\usepackage{natbib}

%%% Sample ICM/MCM Contest Submission
%%%
%%% Based on sample senior thesis document
%%% Last modified by Jeremy Rouse
%%%   Summer 2000
%%%
%%% and on the LaTeX Hints document
%%% created by C.M. Connelly <cmc@math.hmc.edu>
%%%   Copyright 2002-2012


%%% ---------------
%%% Local Command and Environment Definitions

%%% If you have any local command or environment definitions, put them
%%% here or in a separate style file that you load with \usepackage.

% \newtheorem declarations
\newtheorem{Theo1}{Theorem}
\newtheorem{Theo2}{Theorem}[section]
\newtheorem{Lemma}[Theo2]{Lemma}
% Each of the above defines a new theorem environment.
% Multiple theorems can be done in the same environment.
% Theo2's number is defined by the subsection it's in.
% Theo3 uses the same numbering counter and numbering system as
% Theo2 (that's the meaning of [Theo2]).


%%% You probably won't want any of the following commands, which are
%%% here to allow various the names of commands, make examples typeset
%%% properly, and so on.  You can, of course, use them as examples for
%%% your own user-defined commands.
\newcommand{\bslash}{\symbol{'134}}%backslash
\newcommand{\bsl}{{\texttt{\bslash}}}
\newcommand{\com}[1]{\bsl\texttt{#1}\xspace}
\newcommand{\file}[1]{\texttt{#1}\xspace}

\newcommand{\pdftex}{PDF\tex}
\newcommand{\pdflatex}{PDF\latex}
\newcommand{\acronym}[1]{\textsc{#1}\xspace}
\newcommand{\key}[1]{\textsf{\emph{#1}}\xspace}
\newcommand{\class}[1]{\textsf{#1}\xspace}
\newcommand{\package}[1]{\textsf{#1}\xspace}
\newcommand{\env}[1]{\texttt{#1}\xspace}
\newcommand{\prog}[1]{\texttt{#1}\xspace}
\newcommand{\command}[1]{\texttt{\bsl{}#1}\xspace}
\newcommand{\ctt}{\texttt{comp.text.tex}\xspace}
\newcommand{\tex}{\TeX\xspace}
\newcommand{\latex}{\LaTeX\xspace}

%%% Note that the \xspace command comes from the xspace package.  It
%%% allows you type a command that inserts text without having to
%%% worry about how you ``end'' that command.
%%%
%%% Without \xspace, you would need to end a command with a backslash
%%% followed by a space or with an empty set of braces if you followed
%%% the command with a space.  For example,
%%%
%%%   \foo is a very important algorithm.
%%%
%%% might produce 
%%%  
%%%   The foobarbaz algorithmis a very important algorithm.
%%%  
%%% whereas with the \xspace command, the same code would produce
%%%  
%%%   The foobarbaz algorithm is a very important algorithm.
%%%  
%%%  If you need to butt a command that produces text against a letter
%%%  of some sort -- say, to pluralize it -- you need to tell TeX
%%%  where your command name ends so that it expands the correct
%%%  macro.  So you might do
%%% 
%%%    \bar{}s are very busy creatures.


%%% TeX has an amazingly good hyphenation algorithm, but sometimes it
%%% gets confused and needs some help.
%%%
%%% For words that only occur once or twice, you can insert hints
%%% directly into your text, as in
%%%
%%%    our data\-base system is one of the most complex ever devised
%%%
%%% For words that you use a lot, and that seem to keep ending up at
%%% the end of a line, however, inserting the hints each time gets to
%%% be a drag.  You can use the \hyphenation command  to globally tell
%%% TeX where to hyphenate words it can't figure out on its own.

\hyphenation{white-space}



%%% End Local Command and Environment Definitions
%%% ---------------


%%% ---------------
%%% Title Block

\title{\latex Hints for ICM/MCM Contest Reports}

%%% Which contest are you taking part in?  (Just one!)

\contest{ICM/MCM}

%%% The question you answered.  (Again, just the one.)

\question{Report Sample}

%%% Your Contest Team Control Number
\team{12345}

%%% A normal document would specify the author's name (and possibly
%%% their affiliation or other information) in an \author command.
%%% Because the ICM/MCM Contest rules specify that the names of the
%%% team members, their advisor, and their institution should not
%%% appear anywhere in the report, do *not* define an \author command.

%%% Defining the \date command is optional.  If you leave it blank,
%%% your document will include the date that the file is typeset, in
%%% the form  ``Month dd, yyyy''.

% \date{}

%%% End Title Block
%%% ---------------

\begin{document}

%%% ---------------
%%% Summary

\begin{summary}
  This document is meant to give you a quick introduction to \TeX\ and
  \LaTeX.  It covers a lot of material, but still barely manages to
  scratch the surface.  It should provide you with some inspiration
  and, I hope, with some useful code you can copy, modify, and use in
  your report.

  You should use the \file{blank-template.tex} file as a basis for
  your report rather than this file.  Be sure to change its name to
  something sensible (maybe your team control number), and to set the
  values of the \com{title}, \com{question}, and \com{team} commands
  to appropriate values.

  Good luck!

  \hfill{}-- Claire
\end{summary}
 
%%% End Summary
%%% ---------------

%%% ---------------
%%% Print Title Block, Contents, et al.

\maketitle
\tableofcontents

%%% Uncomment the following lines if you have figures or tables in
%%% your report:
\listoffigures
\listoftables  
 
%%% End Print Title Block, Contents, et al.
%%% ---------------



\section{Introduction: What Is \latex?}%
\label{sec:introduction}

\latex is a tool that allows you to concentrate on your writing while
taking advantage of the \tex typesetting system to produce
high-quality typeset documents.

\latex's benefits include
\begin{enumerate}
\item Standardized document classes
\item Structural frameworks for organizing documents
\item Automatic numbering and cross-referencing of structural elements
\item ``Floating'' figures and tables
\item High-level programming interface for accessing \tex's
  typesetting capabilities
\item Access to \latex extensions through loading ``packages''
\end{enumerate}


\section{Structured Writing}%
\label{sec:structured-writing}

Like HTML,\footnote{HyperText Markup Language} \latex is a markup
language rather than a \acronym{Wysiwyg}{}\footnote{What You See Is
  What You Get.}  system.  You write plain text files that use special
\key{commands} and \key{environments} that govern the appearance and
function of parts of your text in your final typeset document.


\subsection{Document Classes}%
\label{sec:document-classes}

The general appearance of your document is determined by your choice
of \key{document class}.  Document classes also load \latex packages
to provide additional functionality.

\latex provides a number of basic classes, including \class{article},
\class{letter}, \class{report}, and \class{book}.  There are also a
large number of other document classes available, including
\class{amsart} and \class{amsbook}, created by the American
Mathematical Society and providing some additional mathematically
useful structures and commands; \class{foils}, \class{prosper}, and
\class{seminar}, which allow you to create ``slides'' for
presentations; the math department's \class{thesis} class, for
formatting senior theses; and many journal- or company-specific
classes that format your document to match the ``house style'' of a
particular periodical or publisher.


\subsection{Packages}%
\label{sec:packages}%
\label{sec:ctan}

\latex packages, or \key{style files}, define additional commands and
environments, or change the way that previously defined commands and
environments work.  By loading packages, you can change the fonts used
in your document, write your document in a non-English language with a
non-\acronym{Ascii} font encoding, include graphics, format program
listings, add custom headers and footers to your document, and much
more.

A typical \tex installation includes hundreds of style files, and
hundreds more are available from the Comprehensive \tex Archive
Network (CTAN), at \url{http://www.ctan.org/}.


\subsection{Structural Commands}%
\label{sec:structural-commands}

\begin{table}
\centering
\begin{tabular}{ll}
\toprule
Command           & Notes                               \\
\midrule
\com{part}        & \class{book} \& \class{report} only \\
\com{chapter}     &\class{book} \& \class{report} only  \\
\com{section}                                           \\
\com{subsection}                                        \\
\com{subsubsection}                                     \\
\com{paragraph}                                         \\
\com{subparagraph}                                      \\
\bottomrule
\end{tabular}
\caption[Structural commands in \latex]{Structural commands in \latex.}%
\label{tab:structural-commands}
\end{table}

\latex provides a set of structural commands for defining sections of
your document, as shown in Table~\ref{tab:structural-commands}.

Note that the argument to structural commands are moving arguments
(see Section~\ref{sec:fragile-commands}) because they can be reused in
the table of contents or in page headers or footers.  Structural
commands can take an optional argument in which you specify nonfragile
commands or a shorter version of the actual section title that fits.
You'll generally know when you need to provide an optional argument by
\TeX's behavior.


\subsection{Labels and References}%
\label{sec:labels-and-references}

Sections are numbered automatically by \latex during typesetting.  If
you change your mind and decide that a subsection should be promoted
to a section, or moved to the end of your document, the sections will
be renumbered so that the numbers are consistent. 

Sections can also be \command{label}{}ed with a tag such as 
\begin{quote}
\begin{verbatim}
\section{Our Complicated Equations}%
\label{sec:complicated-eqs}
\end{verbatim}
\end{quote}
and referred to with a \command{ref} or \command{pageref} command, as
in
\begin{quote}
\begin{verbatim}
In Section~\ref{sec:complicated-eqs}, we pointed out...
\end{verbatim}
\end{quote}
or
\begin{quote}
\begin{verbatim}
On page~\pageref{fig:gordian-knot}, we illustrated...
\end{verbatim}
\end{quote}

\latex substitutes the correct section number when typesetting your
document.

The same commands can be used with numbered environments such as
\env{equation}, \env{theorem}, and so forth.

Use \emph{meaningful} labels---labeling a section as \texttt{sec12}
may seem useful, but it will be confusing if you end up moving it to a
different place in the document and its number changes to Section~34.
It's also easier to remember what reference you want if you use a
meaningful name.

You may also want to impose some additional organization through the
use of \emph{namespaces}, as I've done in this document.  Rather than
give different types of objects undistinguished labels, I precede
section labels with \texttt{sec:}, equations with \texttt{eq:},
figures with \texttt{fig:}, tables with \texttt{tab:}, and so on.

Emacs with Aux\tex and Ref\tex gives you easy access to these labels,
as do many other editors with \tex-specific features.  It's much
easier to find the particular label you're looking for if you have
some additional information to help you.  Adding the prefixes also
reminds you of what text should precede the \com{ref} command.


\subsection{Commands}

\latex uses commands for changes that are very limited in scope (a few
words) or are unlimited in scope (the rest of a document).  For
example, the commands
\begin{quote}
\begin{verbatim}
\textbf{bold} 
\emph{italic (emphasized)} 
\textsf{sans serif}
\end{verbatim}
\end{quote}
produce the following output in a typeset document:
\begin{quote}
\textbf{bold} \emph{italic (emphasized)} \textsf{sans serif}
\end{quote}

These are ``commands with arguments''---the command itself starts with
a backslash (\bsl), and its \key{argument} appears inside braces
{\verb+{ }+).  Some commands may also have \key{optional arguments},
which are typed inside brackets (\verb+[ ]+).

There are also commands that take no arguments, such as
\command{noindent}, \command{raggedright}, and \command{pagebreak}.

You can define your own commands, as discussed in
Section~\ref{sec:customization}.


\subsection{Environments}%
\label{sec:environments}

\latex provides a number of \key{environments} that affect the
appearance of text, and are generally used for more structurally
significant purposes.  For example, the commands listed above are
typeset inside a \env{verbatim} environment typed inside a \env{quote}
environment.  Their results were typeset inside a \env{quote}
environment.

Environments use special commands to start and close---\command{begin}
and \command{end}, followed by the name of the environment in braces,
as in
\begin{quote}
\begin{verbatim}
\begin{quote}
  ``This is disgusting---I can't eat this.  That arugala is so
  bitter\ldots{} It's like my algebra teacher on bread.''
  \flushright -- Julia Roberts in \emph{Full Frontal}
\end{quote}
\end{verbatim}
\end{quote}
producing
\begin{quote}
  ``This is disgusting---I can't eat this.  That arugala is so
  bitter\ldots{} It's like my algebra teacher on bread.''
  \flushright -- Julia Roberts in \emph{Full Frontal}
\end{quote}

Some environments may take additional arguments in braces (required)
or brackets (optional).

Note that the order in which environments nest is extremely important.
If you type an environment inside another environment, the inner
environment must be \command{end}{}ed \emph{before} the second
environment is closed.  It's also vitally important that you have an
\command{end} line for each \command{begin} line, or \latex will
complain.

\subsubsection{The \env{document} Environment and the Preamble}%
\label{sec:document-environment}

The most important environment is the \env{document} environment,
which encloses the \key{body} of your document.  The code before the
\command{begin}\verb+{document}+ line is called the \key{preamble},
and includes the all-powerful \command{documentclass} command, which
loads a particular document class (see
Section~\ref{sec:document-classes}); optional \command{usepackage}
commands, which load in additional \latex packages (see
Section~\ref{sec:packages}); and other setup commands, such as
user-defined commands and environments, counter settings, and so
forth.

I generally also include the commands defining the title, author, and
date in my preambles, but other people include them just after
\command{begin}\verb+{document}+, before the \command{maketitle}
command, which creates the title block of your document.

\subsubsection{Math Environments}%
\label{sec:math-environments}

One of the major hallmarks of \tex is its ability to typeset
mathematical equations.

The two primary ways of doing so are with the use of \key{inline} and
\key{display math environments}.  These environments are used so
often that there are shorthands provided for typing them.  Inline math
environments, such as $a^2 + b^2 = c^2$, can be typed as
\begin{quote}
\begin{verbatim}
\begin{math}
a^{2} + b^{2} = c^{2}
\end{math}
\end{verbatim}
\end{quote}
or
\begin{quote}
\begin{verbatim}
$a^{2} + b^{2} = c^{2}$.
\end{verbatim}

\end{quote}

Display math environments set your equation apart from your running
text.  They're generally used for more complicated expressions, such
as
\[
f(x) = \int \left( \frac{x^2 + x^3}{1} \right)dx
\]
which can be typed as

\begin{quote}
\begin{verbatim}
\begin{displaymath}
f(x) = \int \left( \frac{x^2 + x^3}{1} \right)dx
\end{displaymath}
\end{verbatim}
\end{quote}
or
\begin{quote}
\begin{verbatim}
\[
f(x) = \int \left( \frac{x^2 + x^3}{1} \right)dx
\]
\end{verbatim}
\end{quote}

Generally, you'll want to use the \verb+$+ %$ <- fool font-lock-mode

delimited form for inline math, and the \com{[} \com{]} form for
display math environments.  [Besides being easy to type, these forms
are \key{robust}, which means that they can be used in \key{moving
  arguments}, elements that \tex may need to typeset in more than one
place (such as a table of contents) or adjust (such as footnotes).]

\paragraph{The \env{equation} Environment}%
\label{sec:equation-environment}

You'll probably want to use the \env{equation} environment for any
formula you plan to refer to.  \latex not only typesets the contents
of an \env{equation} environment in display mode, it also numbers it,
as in
\begin{equation}
\label{eq:myequation}
f(x) = \int \left( \frac{x^2 + x^3}{1} \right)dx
\end{equation}
written as
\begin{quote}
\begin{verbatim}
\begin{equation}
\label{eq:myequation}
f(x) = \int \left( \frac{x^2 + x^3}{1} \right)dx
\end{equation}
\end{verbatim}
\end{quote}

Note that you can refer to this formula as
Equation~\ref{eq:myequation} with
\begin{verbatim}
\ref{eq:myequation}.
\end{verbatim}


\subsection{Fonts}%
\label{sec:fonts}

Generally you'll want to let \latex handle the fonts for you---Knuth's
Computer Modern fonts are used by default, and include a wide range of
variations that can cover most any use you can think of.  

If you want to get fancy (and portable; see
Section~\ref{sec:fuzzy-fonts}), you can use Type~1 PostScript fonts,
such as Times, Palatino, Utopia, and so forth.  These font sets are
accessible with packages with names like \package{times},
\package{palatino}, and \package{utopia}.  There are others, as
well---a command such as \com{locate psnfss | grep sty} will find most
of them.

You can also get fonts from CTAN (see Section~\ref{sec:ctan}), both
bitmap and Type 1.  There's even support for TrueType fonts in some
\TeX\ systems.

\subsubsection{Font Commands}%
\label{sec:font-commands}

Most of your concern about fonts is probably related to what you're
writing.  You might want some \emph{emphasized} or \textbf{bold} text
to stress a point or highlight a key term.  Filenames might be set in
\texttt{typewriter text} (although you should consider using the
\package{url} package to help you out---by default, text set in
typewriter text isn't hyphenated, which can lead to some unattractive
line breaks).

You can also set text in \textsf{sans serif} or \textsc{small caps}.
Table~\ref{tab:font-commands} shows you some of the most commonly used
font commands provided by \latex.

\begin{table}[htbp]
  \centering
  \begin{tabular}{ll}
    \toprule
    Command       &   Result\\
    \midrule
    \com{emph}    &   \emph{emphasized text}\\
    \com{textsf}  &   \textsf{sans-serif text}\\
    \com{texttt}  &   \texttt{typewriter text}\\
    \com{textbf}  &   \textbf{bold text}\\
    \com{textsc}  &   \textsc{small caps text}\\
    \com{textsl}  &   \textsl{slanted text}\\
    \com{textit}  &   \textit{italic text}\\
    \bottomrule
  \end{tabular}
  \caption[Commonly used font commands]{Commonly used font commands.}
  \label{tab:font-commands}
\end{table}

I recommend that you use \com{emph} in preference to \com{textit}, and
use \com{textbf} sparingly.  \com{emph} is a smarter command than
\com{textit}---it switches back to the roman font when necessary.  For
example, \emph{She loved \emph{Scooby Doo}.} versus \textit{He loved
  \textit{Titanic}.}

For complicated font changes, or for special font usages that you're
typing a lot, creating a macro (Section~\ref{sec:customization}) is
the way to go.  I often just write, tossing in custom commands as I
go, and waiting to define them until just before I compile the
document.


\subsection{Customization}%
\label{sec:customization}

The main advantage of using commands and environments is that they
allow you to organize your writing.  A useful side-effect is that you
can change your mind about the way an element is typeset, and change
all the appearances of that element in document by editing one piece
of code.  For example, in this document the names of environments have
been set in ``typewriter text'', using a command I created called
\command{env}, which is defined as
\begin{quote}
\begin{verbatim}
\newcommand{\env}[1]{\texttt{#1}\xspace}
\end{verbatim}
\end{quote}

All I have to do to make the names of all the environments in the
document appear in sans-serif type instead is to change that one line
to
\begin{quote}
\begin{verbatim}
\newcommand{\env}[1]{\textsf{#1}\xspace}
\end{verbatim}
\end{quote}

You can do the same with almost anything you can conceptualize---key
terms, people's names (especially names of people from
non-English-speaking countries), files, functions, and so on.


\section{Mathematical Notation}%
\label{sec:mathematical-notation}

As we saw in Section~\ref{sec:math-environments}, math is typed into
one of several kinds of math environments.  Choose your environment
based on the context and importance of the content.  Any formula you
plan to refer to should be typed in an \env{equation} environment (or
a similar environment that supports labels).

You should punctuate your mathematics as if the formulae were normal
parts of English sentences.  Reading them aloud is often a useful
method for ensuring that you have all the commas in the right places.
Where appropriate, you should also follow a displayed formula at the
end of a sentence with a period.

\subsection{Sums and Products}%
\label{sec:sums-n-products}

It's easy to typeset sums and products.  For example,
\begin{equation}
f(n) = \sqrt[n]{\sum_{k=1}^{n} {n \choose k} f \left( n - k \right)},~
\prod_{n=2}^{\infty} \frac{n^{3}-1}{n^{3}+1} = \frac{2}{3}.
\end{equation}

%%% The ~ in the equation puts a nonbreaking space (equivalent to an
%%% interword space in text mode) between the two halves of the equation.
%%%
%%% Also, note that the use of the \choose command here causes the
%%% amsmath package to issue the warning
%%%
%%%    Package amsmath Warning: Foreign command \atopwithdelims;
%%%    (amsmath)                \frac or \genfrac should be used instead
%%%    (amsmath)                 on input line 557.
%%%
%%% amsmath would prefer the use of the \binom command it supplies.


\subsection{Matrices}%
\label{sec:matrices}

It's a little more difficult to create matrices, but not too bad:
%%% In LaTeX, & is the alignment tab, and separates columns. \\ is the end of
%%% line marker, and separates rows. The ccc denotes that there are three
%%% columns.  The array environment and the tabular environment are
%%% more or less identical, so what goes here also applies to a table.
%%%
\begin{equation}
\left[ \begin{array}{ccc}
     2 & 1 & 2\\
     1 & 0 & 2\\
     2 & 1 & 1
     \end{array} \right]
\left[ \begin{array}{ccc}
     -2 & 1 & 2\\
     3 & -2 & -2\\
     1 & 0 & -1
     \end{array} \right] = 
\left[ \begin{array}{ccc}
     1 & 0 & 0\\
     0 & 1 & 0\\
     0 & 0 & 1
     \end{array} \right].
\end{equation}


\subsection{Symbols}%
\label{sec:symbols}

\LaTeX provides an enormous number of symbols.  Additional packages
(loaded with \com{usepackage}) may provide additional symbols and
fonts.

For example, $\mathbb{N}$, $\mathbb{Z}$, $\mathbb{Q}$, $\mathbb{R}$,
and $\mathbb{C}$ require you to load the \package{amsfonts} package
(which is automatically loaded by the \texttt{icmmcm} class).  These
symbols are generated by \com{mathbb}, which only works in math mode.

Subscripts and superscripts are easy---\verb!$a_n$!  produces $a_n$,
and \verb!$x^2$! produces $x^2$.  Ordinal numbers, such as
$3^{\textrm{rd}}$, $n^{\textrm{th}}$, and so forth,\footnote{Some
  fonts may include their own ordinals that can be accessed with
  special commands.} can be produced with code like
\verb!$3^{\textrm{rd}}$!, \verb!$n^{\textrm{th}}$!.
Equation~\ref{eq:superscript} shows a formula with a superscript.
\begin{equation}
\label{eq:superscript}
 \int_{0}^{\pi} \, \cos^{2n+1} x \, {\rm d} x = 0 \qquad
\forall \, n \in \mathbb{N}. 
\end{equation}
Notice that \com{cos} produces a nice roman ``$\cos$'' within math
mode. There are similar commands for common functions like \com{log},
\com{exp}, and so forth.  More can be defined with the
\com{DeclareMathOperator} command provided by the \package{amsmath}
package.

You can stack symbols over other symbols. In math formulas,
\begin{equation}
  m\ddot{x} + \gamma\dot{x} + kx  = 0,
\end{equation}
or to produce diacritical accents, as in
\begin{quote}
  Paul Erd\H{o}s s'est reveill\'{e} t\^{o}t pour enseigner le
  fran\c{c}ais \`{a} son fr\`{e}re et sa s\oe{}ur.
\end{quote}

\LaTeX{} has lots of Greek letters and ellipses too, some of which are
shown in Figure~\ref{fig:greek-symbols}.

\begin{figure}
  \begin{center}
    \begin{equation}
      \sqrt{ 
        \left[ 
          \begin{array}{cccccc}
            \alpha & \beta    & \gamma & \delta  & \epsilon & \zeta  \\
            \eta   & \theta   & \iota  & \kappa  & \lambda  & \mu    \\
            \nu    & \xi      & o      & \rho    & \pi      & \sigma \\ 
            \tau   & \upsilon & \phi   & \chi    & \psi     & \omega \\
            \Gamma & \Delta   & \Theta & \Lambda & \Xi      & \Pi    \\
            \Sigma & \Upsilon & \Phi   & \Psi    & \Omega   & \varphi\\
            \cdots & \ldots   & \vdots & \ddots  & :        & \cdot
          \end{array} 
        \right ] }.
    \end{equation}
  \end{center}
  \caption[Greek letters and some symbols]{Greek letters and some symbols.}%
  \label{fig:greek-symbols}
\end{figure}

See \cite{gratzer-mil}, pp.~455--474, or \cite{kopka-daly-guide},
pp.~123--127, for lists of the symbols available.  Intext, you might
see some of these symbols used as
\begin{quote}
  The Strong Induction Principle asserts that if a statement holds for
  the integers $1$,~$2$,\dots,~$n$, and if whenever it holds for $n =
  1$, \dots,~$k$ then it also holds for $n = k+1$, then the statement
  holds for the integers $1$,~$2$,~$3$, $\ldots\,$ Using this
  Principle, it can be shown that $1+2+\cdots+n = n(n+1)/2$ for all
  positive integers~$n$.
\end{quote}
Notice that in the lists of integers, the ellipsis was made using the
\com{ldots} command, and that the periods were nicely spaced between
the commas. In the sum, the dots were made with \com{cdots} and were
centered on the line.  The \package{amsmath} package provides a
``smart'' \com{dots} command that can generally get things right based
on the context.

So, with \com{dots} alone, the previous examples come out as
\begin{quote}
$1$,~$2$,~\dots,~$n$\\
$n = 1$, \dots,~$k$\\
$1$,~$2$,~$3$, $\dots\,$\\
$1+2+\dots+n = n(n+1)/2$
\end{quote}

The general $n \times n$ matrix can be typeset as follows:
\begin{equation}
\left[
\begin{array}{cccc}
a_{11} & a_{12} & \ldots & a_{1n}\\
a_{21} & a_{22} & \ldots & a_{2n}\\
\vdots & \vdots & \ddots & \vdots\\
a_{n1} & a_{n2} & \ldots & a_{nn}\\
\end{array} 
\right].
\end{equation}


A fine point: lists of numbers that you're using in a mathematical
sense (as opposed to dates, numbers of objects, etc.) should be typed
in math mode.  For example, $341$, $541$, $561$, and $641$.  The same
numbers without math mode are 341, 541, 561, and 641.  Depending on
the fonts and packages that you're using, you may notice a little bit
more space around the first set than the second.  With some packages,
numbers intext may be set using old-style figures by default, as in
\oldstylenums{341}, \oldstylenums{541}, \oldstylenums{561}, and
\oldstylenums{641}.  %%% But without the \oldstylenums commands!


\subsection{More Math}

In Fourier analysis, we talk about the $z$-domain.

If $a$ is an even number, then
\[ a + \phi(a) < \frac{3 a}{2}, \]
and
\[ \sigma(a) > \frac{2^{\alpha+1}-1}{2^{\alpha}} \, a \geq \frac{3
  a}{2}, \]
where $\alpha$ is the greatest power of 2 that divides $a$, $\phi(a)$
is the number of integers less than $a$ and relatively prime
to $a$, and $\sigma(a)$ is the sum of the divisors of $a$ (including
$1$ and $a$). 

Typeset a piecewise function using the \env{cases} environment (from
the \package{amsmath} package) as follows:
\[ |x| = 
\begin{cases} 
   x, & {\rm if~}x \geq 0;\\
  -x, & {\rm otherwise.}\\
\end{cases}
\]

In frosh physics, students come to know the true meaning of
$\mathbf{F} = m \mathbf{a}$, $E = mc^{2}$, and $-\frac{\hbar^{2}}{2m} 
\nabla^{2} \psi + V \psi = i \hbar \frac{\partial \psi}{\partial t}$.

\subsection{Aligning Equations}%
\label{sec:aligning-equations}

In days gone by, people used the \env{eqnarray} environment to align
equations.  \env{eqnarray} has generally been replaced by \env{align}
and some variants such as \env{flalign}, which places the leftmost
column as far left as possible and the rightmost column as far right
as possible; \env{alignat}, which allows you to specify the spacing;
and more.  See \cite{amsmath-doc}, \cite{gratzer-mil}, \cite{lamport},
or \cite{kopka-daly-guide} for more information about the
alternatives.

In Equations~\ref{eq:eqnarray}--\ref{eq:eqnarray-last}, the $=$ signs
have been aligned using the \env{eqnarray} environment.
\begin{eqnarray}{3}%
\label{eq:eqnarray}
   x^4 + \frac{1}{x^4} & = & x^4 + 4 x^2 + 6 + \frac{4}{x^2} + \frac{1}{x^4} -
   4 x^2 - 6 - \frac{4}{x^2}\\
   & = & \left( x + \frac{1}{x} \right)^4 - 4 x^2 - 6 - \frac{4}{x^4}\\
   & = & \left( x + \frac{1}{x} \right)^4 - 4 x^4 - 8 - \frac{4}{x^4} + 8 - 6\\
   & = & \left( x + \frac{1}{x} \right)^4 - 4 \left( x + \frac{1}{x} \right)^2
   - 6.%
\label{eq:eqnarray-last}
\end{eqnarray}

Equations~\ref{eq:align}--\ref{eq:align-last} show the same set
of equations aligned with the \env{align} environment.
\begin{align}{3}%
\label{eq:align}
   x^4 + \frac{1}{x^4} &= x^4 + 4 x^2 + 6 + \frac{4}{x^2} + \frac{1}{x^4} -
   4 x^2 - 6 - \frac{4}{x^2}\\
   &=  \left( x + \frac{1}{x} \right)^4 - 4 x^2 - 6 - \frac{4}{x^4}\\
   &=  \left( x + \frac{1}{x} \right)^4 - 4 x^4 - 8 - \frac{4}{x^4} + 8 - 6\\
   &=  \left( x + \frac{1}{x} \right)^4 - 4 \left( x + \frac{1}{x} \right)^2
   - 6.%
\label{eq:align-last}
\end{align}

\subsection{Adjusting Spacing}%
\label{sec:math-spacing}

Sometimes you need to adjust the spacing and fonts inside integrals.
Typically, the ``d'' (as in $\textrm{d}x$) is set in Roman type.
Rather than
\begin{equation} 
\int \int \frac{1}{1-xy} dx dy.
\end{equation}
you want
\begin{equation}
\int \! \! \! \int \, \frac{1}{1-xy} \, \, {\rm d}x \, {\rm d}y.
\end{equation}
The integral signs have been moved together using the ``negative
space'' command \com{!}.  Extra space has been added between the
elements of integration, $\textrm{d}x$ and $\textrm{d}y$, and between
those elements and the integrand with the ``thin space'' command,
\com{,}.

Table~\ref{tab:math-spacing} shows the different spacing commands
available in math mode.  There are additional spacing commands
provided by the \package{amsmath} package, not shown here.
\begin{table}[htbp]
  \centering
  \begin{tabular}{llll}
    \toprule
            &                &\multicolumn{2}{l}{\LaTeX\ Command}\\
    \cmidrule(l){3-4}
    \multicolumn{2}{l}{Name} &  Short   & Long                   \\
    \midrule
     \multicolumn{2}{l}{Positive Space}                          \\
     \quad{}& thinspace      &  \com{,} & \com{thinspace}        \\
            & medspace       &  \com{:}                          \\
            & thickspace     &  \com{;}                          \\
            & 1 em           &          & \com{quad}             \\
            & 2 em           &          & \com{qquad}            \\
    \addlinespace
    \multicolumn{2}{l}{Negative Space}                           \\
            & thinspace      &  \com{!} & \com{negthinspace}     \\
    \bottomrule
  \end{tabular}
  \caption[\protect\LaTeX\ math spacing commands]{\protect\LaTeX\ math spacing commands.}%
  \label{tab:math-spacing}  
\end{table}


\subsection{Specifying Equation Numbers or Names}%
\label{sec:specifying-eq-numbers}

All the equations you've seen so far are numbered consecutively.  You
can specify a number (or name) for a single equation by placing the
formula in a display math environment (but not an \env{equation}
environment) and giving the desired number or name as the argument to
an \com{eqno} command.  For example,
\[ \int_{0}^{\infty} \! e^{-x^{2}} \, {\rm d}x =
\frac{\sqrt{\pi}}{2}. \eqno (42), \]
or
\[ \int_{0}^{\infty} \! e^{-x^{2}} \, {\rm d}x =
\frac{\sqrt{\pi}}{2}. \eqno (\textrm{cool formula}), \]

Note that you have to specify parentheses in the argument to the
\com{eqno} command.  If you name a formula, you also have to enclose
the text within a command such as \com{mathrm}, or it will be set as
if it was a string of variables (and without any spaces).
 
If you'd like to have many aligned equations without numbers, use the
starred form of the \env{align} environment, \env{align*}, as in
\begin{quote}
The area $K$ of $\triangle ABC$ is given by
\begin{align*}
  K &= \frac{a h_{a}}{2}\\
    &= \frac{a b}{2} \sin C\\
    &= rs\\
    &= \sqrt{s (s-a) (s-b) (s-c)}\\
    &= \frac{a b c}{4 R}\\
    &= \frac{a^{2} \sin B \sin C}{2 \sin A}\\
    &= 2 R^{2} \sin A \sin B \sin C,
\end{align*}
where $A$, $B$, and $C$ are the angles in $\triangle ABC$, $r$ is the
radius of the inscribed circle, $R$ is the radius of the circumscribed
circle, $s$ is one-half of the perimeter, and $h_{a}$ is the length of
the altitude from the vertex $A$ to the side $BC$.
\end{quote}

\subsection{Sizing Delimiters}%
\label{sec:sizing-delimiters}

The \com{left} and \com{right} commands, followed by a bracket, brace,
or parenthesis, tell \TeX\ to adjust the size of the delimiter to its
contents.
\begin{equation}
  f(x) = \left(1 + \left(1 + x \right)^{2} \right)^{n}.
\end{equation}
  You can also use commands such as \com{big}, \com{Big},
\com{bigg}, or \com{Bigg} to specify larger delimiters (useful if
  you have multiple levels of delimiters), as in
\[ \big( \quad \Big( \quad \bigg(  \quad \Bigg( \]
or
\[ \Bigg( \bigg( \Big( \big( \, ( a + b ) + c \big) + d \Big) + e \bigg) + f \Bigg) \]


\subsection{Theorems}%
\label{sec:theorems}

The \env{theorem} environment is an easy way to typeset theorems in
the text.  To use it, type a \com{newtheorem} command in the preamble
of your document like
\begin{verbatim}
   \newtheorem{Theo1}{Theorem}
\end{verbatim}

You can then type a theorem using your theorem environment. 

This document includes three such definitions, 
\begin{verbatim}
   \newtheorem{Theo1}{Theorem}
   \newtheorem{Theo2}{Theorem}[section]
   \newtheorem{Lemma}[Theo2]{Lemma}
\end{verbatim}
which show you some of the possibilities available.  Examples of each
appear below.

\subsubsection{A \env{Theo1} Environment}

\begin{Theo1}
The equation $x^4 + y^4 = z^4$ has no solutions where $x$, $y$, and $z$ are
positive integers.
\end{Theo1}

\subsubsection{A \env{Theo2} Environment}

\begin{Theo2}[Wilson]
% Wilson is the title of this theorem
A positive integer $p$ is prime if and only if 
$$(p-1)! \equiv -1 \pmod{p}.$$
\end{Theo2}

\subsubsection{A \env{Lemma} Environment}

\begin{Lemma}
\label{mug}
Prof.~Bernoff's Putnam mug is a multiply connected 2-manifold of genus 1.
\end{Lemma}

\subsection{Proofs}%
\label{sec:proofs}

Adding a proof is very simple. Two positive integers $a$ and $b$ are
amicable if $\sigma(a) = \sigma(b) = a + b$, where $\sigma(N)$ denotes
the sum of the divisors of $N$, as above.  The following is a theorem
with an associated proof.

\begin{Theo2}
There do not exist two consecutive integers which are amicable.
\end{Theo2}
{\bf Proof:}
Since even numbers are annoying, no integers are amicable with even numbers.
Thus, if two consecutive integers are amicable, they are both odd. However,
two consecutive odd numbers do not exist. \hfill $\Box$

\medskip

To create the end-of-proof marker shown here, use 
\verb!\hfill$\Box$!. The \verb!\hfill! makes the box print flush right
at the end of the line, as here.

You can also use the \env{proof} environment provided by the
\package{amsthm} package.  


\section{Figures and Tables---\protect\LaTeX's Float Environments}%
\label{sec:figs-and-tabs}

\LaTeX\ provides two ``float'' environments, \env{figure} and
\env{table}.  Float environments are so called because they can be
typeset on a later page in your document than their location in the
source code.

The \env{table} environment is generally used
for---surprise!---tables.  The \env{figure} environment is often used
for graphs or diagrams, but could also be used for other illustrative
graphics.

The basic float environments don't format their contents specially.
If you want an illustration or table to appear centered, you will need
to type it inside a \env{center} environment or add a \com{centering}
command.

\subsection{Captions}%
\label{sec:captions}

By adding a \com{caption} command, you can specify a caption that will
appear with the float.  Its position depends on where in the
environment you type it---if the command is at the top of the
environment, the caption will be typeset above its contents; if at the
bottom, the caption will appear beneath its contents.  Captions are
usually set at the bottom of a float, but if a particular publisher or
journal prefers the captions on top, you can accommodate them.

Captions should generally be written as brief, complete sentences,
ending with a period.  They should either be capitalized as normal
sentences or use headline capitaliztion---capitalized as you would the
title of a document or a section.  So
\begin{quote}
  Production Statistics from Soviet Russia, 1977--1987.
\end{quote}
or
\begin{quote}
  Production statistics from Soviet Russia, 1977--1987.
\end{quote}
rather than
\begin{quote}
  production statistics from Soviet Russia, 1977--1987
\end{quote}
Whichever style you choose, be consistent!

Avoid explaining the whole float in the caption.  Do your explanation
in the text that refers to the float.  

The \com{caption} command takes an optional argument, which is typed
inside brackets (\verb|[ ]|).  This argument is used in the list of
tables or list of figures in place of your actual caption.


\subsubsection{Fragile Commands and Moving Arguments}%
\label{sec:fragile-commands}

Both arguments to the \com{caption} command are \emph{moving
  arguments} (because \TeX\ can move them).  Some commands are
\emph{fragile}, that is, they produce output that can cause problems
if the typeset text is moved somewhere other than the place that \TeX\ 
originally thought it would be typeset.
  
  To prevent fragile commands from being expanded too early and
  causing problems, you can use the \com{protect} command just before
  the command you want to keep unexpanded.


\subsubsection{Labels}%
\label{sec:labeling-floats}

The \com{label} command for a float is generally typed immediately
after the \com{caption} command.


\subsection{Figures}

\subsubsection{Including Graphics}%
\label{sec:including-graphics}

The \env{figure} environment is often used for including graphic
images.  The state-of-the-art method requires you to load the
\package{graphics} or \package{graphicx} package.  Both packages
provide the same functionality, but take arguments in a slightly
different format.\footnote{The \package{graphicx} package defines
  commands that take their arguments in key--value pairs.}  More
information about the \texttt{graphics} package is available in its
manual, \texttt{grfguide} \cite{carlisle-grfguide}, which is included
in DVI, PostScript, or PDF format with most \TeX\ systems.

The standard graphic format used with \TeX\ is called
\emph{Encapsulated PostScript}, or EPS.  EPS files are special
PostScript files that define their ``bounding box'', may include a
bitmap representation for use in previewers, and are restricted from
using some PostScript operators.

EPS files are generally created with a vector graphics application
such as Adobe Illustrator, Dia, OmniGraffle, or Visio.  They can also
be created from \TeX\ files using the \texttt{-E} flag with
\prog{dvips}.

With the development of PDF\TeX, generating Portable Document Format
files has become much easier.  PDF\TeX\ requires that your graphics
are also PDF files (or PNGs, if they're bitmaps).  See
Section~\ref{sec:pdftex-graphics} for some hints.

Both formats use the same command, however: \com{includegraphics}.

The following code produces the graphic in
Figure~\ref{fig:an-eps-graphic}:
\begin{verbatim}
   \begin{figure}[ht]
     \begin{center}
       \scalebox{.50}{\includegraphics{shapes}}
     \end{center}
     \caption[Some shapes]{Some shapes.}%
     \label{fig:an-eps-graphic}
   \end{figure}
\end{verbatim}
\begin{figure}[ht]
  \begin{center}
    \scalebox{.50}{\includegraphics{shapes}}
  \end{center}
  \caption[Some shapes]{Some shapes.}%
  \label{fig:an-eps-graphic}
\end{figure}

Notice that we didn't specify the \texttt{.eps} extension in the
filename argument to the \com{includegraphics} command.  By dropping
the extension, we can typeset this document with PDF\LaTeX\ by simply
changing the driver argument in the line that loads the
\package{graphics} package (provided that we have a PDF version of the
graphic, of course).  The \com{includegraphics} package searches for
different graphic formats depending on the driver it's using.


\subsubsection{\LaTeX\ Diagrams}

\LaTeX\ can also be used to create both simple pictures and
sophisticated diagrams.  Figure~\ref{fig:step-function} shows a graph
created with the \env{picture} environment.

\begin{figure}[ht]
\begin{center}
\begin{picture}(256,300)(0,0)
\put(0,0){\line(1,0){256}}
\put(0,0){\line(0,1){300}}
\put(0,58.6){\line(1,0){128}}
\put(192,82.8){\line(1,0){64}}
\put(160,117.2){\line(1,0){32}}
\put(128,165.7){\line(1,0){16}}
\put(144,234.3){\line(1,0){8}}
\put(128,-2){\line(0,1){4}}
\put(192,-2){\line(0,1){4}}
\put(160,-2){\line(0,1){4}}
\put(144,-2){\line(0,1){4}}
\put(152,-2){\line(0,1){4}}
\put(-2,100){\line(1,0){4}}
\put(-2,200){\line(1,0){4}}
\multiput(154,0)(0,8){35}{\line(0,1){2}}
\put(153,-10){$t$}
\put(-8,-10){$0$}
\put(256,-10){$1$}
\put(-10,100){$1$}
\put(-10,200){$2$}
\put(60,10){$\mathcal{F}_{0}$}
\put(220,10){$\mathcal{F}_{1}$}
\put(172,10){$\mathcal{F}_{2}$}
\put(130,10){$\mathcal{F}_{3}$}
\end{picture}
\caption[A step function]{A step function with a peak at $t$. }%
\label{fig:step-function}
\end{center}
\end{figure}

Chapter~6 of \cite{kopka-daly-guide} describes how you can create diagrams
such as that shown in Figure~\ref{fig:step-function}.


\subsection{Tables}%
\label{sec:tables}

Tables are a complicated subject, not because they're difficult to do
in \latex, but because they're difficult to do \emph{right}.  Most
books on \latex cover tables, but present what Simon Fear, author of
the \package{booktabs} package, calls ``tableaux''.  One such tableau
is illustrated in Table~\ref{tab:tableau}.

\begin{table}[htbp]
  \begin{center}
    \begin{tabular}{||l|lr||} \hline
      gnats     & gram      & \$13.65 \\ \cline{2-3}
                & each      & .01     \\ \hline
      gnu       & stuffed   & 92.50   \\ \cline{1-1} \cline{3-3}
      emu       &           & 33.33   \\ \hline
      armadillo & frozen    & 8.99    \\ \hline
    \end{tabular}
  \end{center}
  \caption[A tableau]{A tableau.  (Taken from \cite{lamport}, pg.~64.)}%
  \label{tab:tableau}
\end{table}

\begin{table}[htbp]
  \begin{center}
    \begin{tabular}{@{}llr@{}} \toprule
      \multicolumn{2}{c}{Item} \\ \cmidrule(r){1-2}
      Animal    & Description & Price (\$)\\ \midrule    
      Gnat      & per gram    & 13.65 \\                 
                & each        & 0.01 \\                  
      Gnu       & stuffed     & 92.50 \\                 
      Emu       & stuffed     & 33.33 \\                 
      Armadillo & frozen      & 8.99 \\ \bottomrule
    \end{tabular}
  \end{center}
  \caption[The tableau as a table]{The tableau as a table.}%
  \label{tab:tableau-table}
\end{table}

Simon argues that such a tableau would be better presented as the
table shown in Table~\ref{tab:tableau-table}.  \emph{The Chicago
  Manual of Style}, \cite{chicago}, and Edward Tufte
\citeyearpar{tufte-vdq} support his assertion, and provide excellent
references and inspiration.

The \package{booktabs} package, which is automatically loaded by the
\file{icmmcm} class, has some special commands for creating lines of
different thicknesses for use as top, bottom, and midrules.  It also
has some code that provides the \com{cmidrule} command, for creating
spanner rules for decked spanner heads.  The rest is up to you and
your style guide.

As an example of a table to strive towards,
Table~\ref{tab:chicago-table} was taken from an example table in
\emph{The Chicago Manual of Style}.  All the tables (except for those
in the section describing tables, alas) in Gr\"{a}tzer's \emph{Math
  into \LaTeX} were prepared with \package{booktabs}.

\begin{table}
\begin{center}
{\hspace{-1in}
\begin{minipage}{\textwidth}
\fontsize{10}{12}\selectfont
\begin{tabular}[c]{lrrrrrr}
\toprule
              & \multicolumn{2}{c}{1900} & \multicolumn{2}{c}{1906} & \multicolumn{2}{c}{1910}\\
\cmidrule(r){2-3}\cmidrule(lr){4-5}\cmidrule(l){6-7}
Party         & \% of Vote  & Seats Won  & \% of Vote  & Seats Won  & \% of Vote  & Seats Won \\
\midrule
\addlinespace
              & \multicolumn{6}{c}{Provincial Assembly}\\
\cmidrule{2-7}
Conservative  & 35.6        &  47        & 26.0        & 37         & 30.9        & 52\\
Socialist     & 12.4        &  18        & 27.1        & 44         & 24.8        & 39\\
Christian Democrat & 49.2   &  85        & 41.2        & 68         & 39.2        & 59\\
Other         & 2.8         &  0         & 5.7         & 1          & 5.1         & 0\\
\addlinespace
Total& 100.0       &  150       & 100.0       & 150        & 100.0       & 150\\
\addlinespace
              & \multicolumn{6}{c}{National Assembly}\\
\cmidrule{2-7}
Conservative  & 32.6        &   4        & 23.8        &  3         & 28.3        & 3\\
Socialist     & 13.5        &   1        & 27.3        &  3         & 24.1        & 2\\
Christian Democrat & 52.0   &   7        & 42.8        &  6         & 46.4        & 8\\
Other         & 1.8         &   0        & 6.1         &  0         & 1.2         & 0\\
\addlinespace
Total& 100.0       &  12        & 100.0       & 12         & 100.0       & 13\\
\bottomrule
\end{tabular}
\end{minipage}
}
\caption[Elections in G\"{o}tefrith province, 1900--1910]{Elections in
  G\"{o}tefrith province, 1900--1910.  (Taken from \cite{chicago},
  pg.~414.)}%
\label{tab:chicago-table}
\end{center}
\end{table}



\section{Typesetting}%
\label{sec:typesetting}

So you've got a \latex source document.  How do you get a typeset
document that you can print or put on the web?

Typesetting a document is referred to as ``\tex{}ing'', ``compiling'',
or ``typesetting''.  Generally, you want to create a PostScript file
(for printing) or a PDF file (for placing on the web).  There are
multiple ways to do both tasks.


\subsection{Getting to Paper}

Starting with a \latex document, \file{foo.tex}, you can create a
PostScript file by running the following commands:
\begin{quote}
\begin{verbatim}
unix% latex foo
unix% dvips -o foo.ps foo
unix% lpr foo.ps
\end{verbatim}
\end{quote}
(On some systems, \prog{dvips} automatically prints your document to
the default printer.  You need to specify the \prog{-o} flag to get a
PostScript file on such systems.)


\subsection{PDF for the Web}

Starting with \file{foo.tex}, you can create a PostScript file by
typing
\begin{quote}
\begin{verbatim}
unix% pdflatex foo
\end{verbatim}
\end{quote}
or with the following sequence of commands:
\begin{quote}
\begin{verbatim}
unix% latex foo
unix% dvips -Ppdf foo -o foo.ps
unix% ps2pdf foo.ps foo.pdf
\end{verbatim}
\end{quote}
(The last step can be replaced by running the PostScript file through
Adobe Acrobat Distiller.)


\subsection{General Comments}

\latex does its numbering (and some other functions) by writing
information to an \key{auxiliary file}.  It then reads that
information in on the next pass, and uses it to typeset references.
Thus you have to run \prog{latex} or \prog{pdflatex} at least twice
whenever you make a change that affects the numbering of elements or
the flow of text across pages.  It's generally good practice to run
\latex three times, or until it stops warning you about possible
changes.

\subsection{Additional Programs}

There are some additional functions, such as indexing and
bibliographies, that use external programs to read auxiliary files and
produce \latex code for inclusion on later runs.  We won't cover those
programs in this document.


\section{Tips and Tricks}

\latex is a very complicated and powerful language.  As a result,
there are many sneaky aspects to it that will cause you problems if
you don't know about them.  Here are a few.

\subsection{Special Characters}%
\label{sec:special-chars}

\tex and \latex have a number of ``special'' characters that are
reserved for use by the language.  Using these characters in your
writing requires you to do a bit of extra work, as shown in
Table~\ref{tab:special-chars}.

\begin{table}
\begin{tabular}{llll}
\toprule
\multicolumn{2}{l}{Character}   & Function     & To Typeset\\
\midrule
\#          & octothorp      & Macro parameter character  & \verb+\#+\\
\$          & dollar sign    & Start/end inline math mode & \verb+\$+\\
\%          & percent sign   & Comment character          & \verb+\%+\\
\&          & ampersand      & Column separator           & \verb+\&+\\
\_          & underscore     & Subscripts, as in $x_2$    & \verb+\_+\\
\{, \}      & braces         & Parameters                 & \verb+\{, \}+\\
\~{}        & tilde          & Nonbreaking space          & \verb+\~{}+ \\
\^{}        & caret          & Superscripts, as in $x^2$  & \verb+\^{}+\\
\verb+\+    & backslash      & Starts commands   & \command{verb}\verb+|\|+ \\
\bottomrule
\end{tabular}
\caption[Special characters in \latex]{Special characters in \latex.}%
\label{tab:special-chars}
\end{table}


\subsection{Comments and Spacing}

You can add comments to your source file that won't appear in your
typeset document by starting them with a \verb+%+.  Any line that
starts with a \verb+%+ will be ``commented out'', and won't be
interpreted.  You can also add a \verb+%+ at the end of a line, with
or without text, and it will make the end of the line disappear.

For example,
\begin{quote}
\begin{verbatim}
% This is a comment line.
This is not a comment line.

This line has a comment at the end%
% This line should be invisible.
of the line.
\end{verbatim}
\end{quote}
will typeset as
\begin{quote}
% This is a comment line.
This is not a comment line.

This line has a comment at the end%
% This line should be invisible.
of the line.
\end{quote}

Notice the lack of a space in ``endof'' on the last line of the
typeset output.  \tex expects a carriage-return character at the end
of a line, and interprets that carriage return as an interword space.
If you comment out the end of a line, you also comment out the
carriage return on that line, and you'll have words run into one
another unless you have a space before the \verb+%+.

\tex collapses multiple spaces into one, and ignores whitespace at the
beginning of a line.  Thus
\begin{quote}
\begin{verbatim}
No spaces.

    Five spaces.

    A tab.
\end{verbatim}
\end{quote}
typesets as
\begin{quote}
No spaces.

    Five spaces.

    A tab.
\end{quote}

(The lines are indented because they are at the start of a paragraph.
You can suppress paragraph indentation with \command{noindent}.)

Paragraphs are delimited by two carriage returns (with or without
whitespace between them).


\subsection{Quotes and Dashes}

Because \tex was designed to do high-quality typesetting, it cares
about which quotation mark and dash you're using, and requires you to
specify the correct punctuation (although most text editors with
special \tex modes will do the substitution for you).

Open and close double quotes---`` and ''---are created by typing
\verb+``+ and \verb+''+, respectively.  The double-quote mark,
\verb+"+, is typeset as " (and is useful for abbreviating ``inches'',
as in 36").

Single-quotes, ` and ', are typed with \verb+`+ and \verb+'+.

There are three basic forms of dashes:
\begin{enumerate}
\item The hyphen, -, is typed as a single dash, \verb+-+
\item The en dash, --, is typed as two dashes, \verb+--+
\item The em dash, ---, is typed as three dashes, \verb+---+
\end{enumerate}

Hyphens are used in hyphenated words, as in ``complex-typesetting
mechanism''.  En dashes are used to indicate ranges, as in ``there are
35--50 of them''.  Em dashes are used to separate independent phrases,
as in ``John believed---honestly believed---that he was right.''

Note that you shouldn't type spaces around any of these dashes---they
run directly against the words on either side, as in \verb+35--50+.


\subsection{Controlling Pagination}%
\label{sec:pagination}

Sometimes you may need to override \LaTeX's choices for line or page
breaks. your document into lines and/or pages.  The \com{pagebreak}
command causes \LaTeX\ to start a new page immediately after the
command appears.  The \verb|\\| command can be used to tell \LaTeX\
where to break a line.

In general, you should let \LaTeX\ have its way, especially if your
document is going to be published by someone else, as they will
undoubtedly have many changes that will have to be made before your
document works for them.  If you do need to tinker with your
document's layout, you should avoid doing so until you're very nearly
done.  If you go back and add or remove text after forcing \LaTeX\ to
do your will, you may find that new blank spaces appear as a result of
your changes.

George Gr\"{a}tzer's \emph{Math into \LaTeX} includes a chapter on
preparing books that covers this topic in depth.


\subsection{Using Graphics with \pdftex}%
\label{sec:pdftex-graphics}

\pdftex supports PDF and JPEG as native graphic file formats.  EPS is
not directly supported---to use EPS figures with \pdftex, you must
first convert your EPS files to PDF.  

If you're using a graphics program such as Adobe Illustrator to
prepare your figures, just save them as PDF instead of (or in addition
to) EPS.

If you don't have access to the tool you used to create your images,
but you still need to convert them, you can use the program
\prog{epstopdf}.

\prog{epstopdf} writes to standard output by default, so you'll have
to redirect the output to a file, as in
\begin{quote}
\begin{verbatim}
unix% epstopdf foo.eps > foo.pdf
\end{verbatim}
\end{quote}

To convert a whole slew of files, you could use a command such as the
following (with the csh):

\begin{quote}
\begin{verbatim}
unix% foreach f ( `find . -type f -name '*.eps'`)
foreach? eps2pdf $f -o=$f:r.pdf
foreach? end
\end{verbatim}
\end{quote}


\subsection{Fonts Look Fuzzy in PostScript or PDF Files}%
\label{sec:fuzzy-fonts}

When Knuth wrote \tex, typesetting was done by trained typesetters
using expensive equipment to cast molten lead into runs of type.
Knuth created his own font family, Computer Modern, by writing a tool
called \MF{}.  \MF{} reads in programs that define various aspects of
every character in a font, and generates bitmap representations of
those characters at a particular resolution, ready for printing.

Unfortunately, bitmaps with resolutions suited for printing look
terrible on screen.  The solution is to use Type 1 PostScript fonts
instead of bitmaps. If you're using \pdftex (or \pdflatex), you get
Type 1 fonts without having to do anything special (but see
Section~\ref{sec:pdftex-graphics}).

If you're using \prog{dvips} to get PostScript as an intermediate step
(using \prog{ps2pdf} or Acrobat Distiller to get PDF), you can force
\prog{dvips} to use Type 1 fonts by specifying the \prog{-Ppdf} flag,
as in
\begin{quote}
\begin{verbatim}
unix% dvips -Ppdf foo.dvi -o foo.ps
\end{verbatim}
\end{quote}


\subsection{Debugging}

One of the trickiest things about using \latex is interpreting
\latex's sometimes cryptic error messages.

In particular, the line numbers that \latex reports are often not the
line numbers where the problem \emph{is}, but the line numbers where
\latex noticed there was a problem.

One useful way of getting a bit more context to help you understand
the problem is to put the line
\begin{quote}
\begin{verbatim}
\setcounter{errorcontextlines}{1000}
\end{verbatim}
\end{quote}
in the preamble of your document, which will provide you with a
(perhaps excessive) amount of context for an error.

The most common errors are probably
\begin{itemize}
\item Using one of the special characters (see
Section~\ref{sec:special-chars})
\item Leaving off or mismatching a brace or bracket
\item Leaving out or swapping arguments to a command or environment
\end{itemize}

If you've tried everything and you can't find the source of an error
message, try the following procedure:
\begin{enumerate}
\item Create a new file, copying your preamble into it
\item Try typesetting it---if you have an error, the problem is in
  your preamble
\item If it typesets, copy half of your document's body into the new
  file, and typeset that
\item If you see your error, then continue halving the document until
  you narrow it down to the problem section
\item If you don't see your error, try the other half
\end{enumerate}



\section{Resources}

There are lots of great resources available for using \tex and \latex.
Here are a few (there are also links available online at
\url{http://www.math.hmc.edu/computing/support/tex/}).


\subsection{Online Documentation}

Much of the documentation for \tex and \latex is available online, as
part of the \tex system.  te\tex, the \tex system installed on the
math lab computers, includes a script called \prog{texdoc} to access
this documentation.  All you have to do is type \prog{texdoc} followed
by a string that you believe is the name of the document you're
looking for.  For example, \prog{texdoc booktabs} will give you the
documentation for the \package{booktabs} package that I used to create
the tables in this document.

Unfortunately, \prog{texdoc} only works for documentation that is
sensibly named.  The authors of the \package{graphics} package, for
instance, called their manual \file{grfguide}.  Still others decided
that \file{manual} was a good name for their manual (after all, it's
the only \emph{manual} in their distribution).

Sometimes you can find documentation using the \prog{locate} command,
which lists all the files on your system that match a string that you
provide.  For example, you could find \file{grfguide} by trying
\prog{locate graphics} and \prog{grep}ping out the results with
\file{texmf} in them, and passing that list to another \prog{grep} for
the string \file{doc}:
\begin{quote}
\begin{verbatim}
unix% locate graphics | grep texmf | grep doc
\end{verbatim}
\end{quote}

Another hard to find, but very useful, document, is the ``User's Guide
for the \package{amsmath} Package'' \citeyearpar{amsmath-doc}, which
is called \file{amsldoc}.


\subsection{UK-TUG FAQ}

The primary list of frequently asked questions in the \tex world is
the UK TUG FAQ, available at
\url{http://www.tex.ac.uk/cgi-bin/texfaq2html}.  If you're not sure
how to do something, or you've got a problem that you're pretty sure
isn't being caused by a typo, check here first.


\subsection{\ctt}

If you can't find an answer in the UK-TUG FAQ, then your next step is
to check \ctt, the Usenet newsgroup devoted to \tex and \latex.
Chances are, whatever your problem is, someone else already had it,
asked about it on \texttt{c.t.t}, and got an answer.  Thanks to Google, Usenet's
past is preserved in an easily searchable format.  Go to Google Groups
(\url{http://groups.google.com/}), type in some search terms, and
check out the answers.  (If you specify \texttt{group:comp.text.tex}
at the end of your search terms, you'll only see results from \ctt.)


\section{Books}

\url{http://www.math.hmc.edu/computing/support/tex/} has some brief
reviews of a number of significant books about \tex and \latex.

My pick for the best introductory/reference book is the third edition
of George Gr\"{a}tzer's \emph{Math into \latex{}}
\citeyearpar{gratzer-mil}.\footnote{Which I edited.}  It's the only book I'm
aware of that discusses the latest version of AMS\latex in depth.  It
also has excellent reference tables and a thorough index.

Another book I highly recommend is Lyn Dupr\'{e}'s \emph{BUGS in
  Writing} \citeyearpar{dupre-bugs}.  Dupr\'{e} is one of Addison Wesley's
senior editors, and has edited many of the most significant books
published by Addison Wesley.  \emph{BUGS} is an accessible guide to
writing clearly and effectively.  It's the kind of book you leave in
the bathroom so you'll always have something interesting and amusing
to read.  Learning how to write better is almost a byproduct!

If you get serious about typesetting, and want to start doing some
fancy page design or want to be sure you're using the right kind of
type, Robert Bringhurst's \emph{The Elements of Typographic Style}
\citeyearpar{bringhurst-elements} will show you the way.

\section{Acknowledgments}

Thanks to Darryl Yong, who wrote a similar document in November, 2000,
from which I borrowed some ideas.  Thanks also to Zeke Burgess, who
wrote the \file{thesis.cls} class and, presumably, the original
version of the \file{samplethesis.tex} file from which I stole some of
the mathematical examples, and Jeremy Rouse, who modified that sample
file.


%%% ---------------
%%% Bibliography

\nocite{*}   %%% Include everything in the thesis.bib file.  Be
             %%% careful---some journals and fields expect you to only
             %%% include references for materials that you have
             %%% actually cited in your paper, others allow you to
             %%% include materials you used as ``background'' without
             %%% actually citing specific pages or passages.
\bibliographystyle{apalike}
\bibliography{icmmcm}

\end{document}



